\documentclass[12pt, a4paper]{article}
\usepackage[utf8]{inputenc}
\usepackage[T1]{fontenc}
\usepackage{amsmath,amssymb}
\usepackage[margin=2.5cm]{geometry}
\usepackage{enumitem}
\usepackage{xcolor}
\usepackage{fancyhdr}
\usepackage{lastpage}

\definecolor{titleblue}{HTML}{1F4E79}
\definecolor{lightblue}{HTML}{EAF3F8}

\pagestyle{fancy}
\fancyhf{}
\lhead{Algorithms Study Plan}
\rhead{Week 1 -- Day 3}
\cfoot{\thepage\ of \pageref{LastPage}}
\setlength{\headheight}{15pt}

\setlength{\parindent}{0pt}
\setlength{\parskip}{0.7em}

\begin{document}

\begin{center}
    {\LARGE\color{titleblue}\textbf{Week 1 -- Day 3}}\\[0.3em]
    {\LARGE\textbf{Exponents and Logarithms}}
\end{center}

\section*{Objective}

Review exponents and logarithms and understand why algorithms that repeatedly divide their input by two usually have 
logarithmic time complexity.

\section*{Exponent Rules}
For a nonzero base $a$, the main exponent rules are: 

\begin{enumerate}[label=\textbf{\arabic*.}, leftmargin=*]
    \item \textbf{Product rule:}
    \[
        a^m \cdot a^n = a^{m+n}
    \]

    \item \textbf{Quotient rule:}
    \[
        \frac{a^m}{a^n} = a^{m-n}
    \]

    \item \textbf{Power rule:}
    \[
        \left(a^m\right)^n = a^{mn}
    \]

    \item \textbf{Zero exponent:}
    \[
        a^0 = 1
    \]
    
    \item \textbf{Negative exponent:}
    \[
        a^{-n} = \frac{1}{a^n}
    \]

\end{enumerate}

\section*{What Does a Logarithm Represent?}

A logarithm answers the following question:

\begin{quote}
    To what exponent must the base be raised to obtain a given number?
\end{quote}

In general, 

\[
    \log_b(x) = y \quad\Longleftrightarrow\quad b^y = x
\]

For example:

\begin{itemize}[leftmargin=*]
    \item $\log_2(8)=3$ because $2^3=8$.
    \item $\log_2(32)=5$ because $2^5=32$.
    \item $\log_2(1024)=10$ because $2^{10}=1024$
\end{itemize}

\section*{Connection with Binary Search}

Starting with $32$ elements and repeatdly dividing the search space by two gives: 

\[
    32 \;\longrightarrow\; 16 \;\longrightarrow\; 8
    \;\longrightarrow\; 4 \;\longrightarrow\; 2 \;\longrightarrow\; 1
\]

It takes five divisions to reach one element, and: 

\[
    \log_2(32)=5
\]

For $1024$ elements, it takes ten divisions because:

\[
    \log_2(1024)=10
\]


Binary Search discards approximately half of the remaining elements during every iterations. Therefore,
the number of iterations grows logarithmically as the input grows.

\begin{center}
    \colorbox{lightblue}{\parbox{0.75\textwidth}{
        \centering
        \textbf{Time complexity:} \(O(\log n)\)
    }}

\end{center}

\end{document}